% 本文件由 LaTeX 排版 Agent 根据有效 translation.json 自动生成。
% author=Tertullian; work_id=0325; title=论忍耐

\chapter{论忍耐}

% structure: p0001 source=h1 target=omitted reason=page-title-already-rendered-by-parent

% structure: p0002 source=h2 target=omitted reason=duplicate-page-title

我向主天主完全承认，我竟敢撰写一篇论忍耐的论文，实在是过于冒失，甚至可以说是厚颜无耻，因为我在实践忍耐方面全然不配——我是一个毫无善德的人。而那些着手阐述和推荐某事物的人，本应首先在该事物的实践上显为出众，并应以自身行为的坚定来调节他们规劝的恒心，以免他们的话语因行为的缺失而羞愧。但愿这种\Quote{羞愧}能带来补救，那样，因未能表现出我们准备提示他人的东西而产生的羞耻，应成为引导我们表现出它的教训；只是某些善事（正如某些恶事一样）的伟大是难以承受的，唯有神圣灵感的恩宠才能真正获得并实践它们。因为至善之物最属于天主；除了拥有它的那一位，没有人能按照祂认为合宜的方式分施它。因此，讨论自己并未被赋予享受的事物，将是一种慰藉；仿效病人的方式，他们既然没有健康，便不知如何缄默不语其福惠。所以我，最可怜的人，常常因不耐心而发热生病，必然要叹息、呼求并坚持不懈地恳求我所没有的忍耐之健康；同时我回忆并在我自己软弱的默想中领会一个真理：信德之健康与主之训导的健全，若非忍耐常伴左右，不易临于任何人。忍耐如此统御着天主之事，以至于若一个人远离忍耐，便不能服从任何诫命，也不能完成任何令主悦纳的工作。忍耐之善，即便那些活在忍耐之外的人，也以最高德行的名号尊崇它。哲学家们（被认为具有相当智慧的生灵）为忍耐安排了如此崇高的地位，以致虽然他们因各派不同的幻想和相互竞争的见解而彼此不和，但唯独对忍耐怀有共同的尊重，他们联合起来将和平赋予这一追求：他们为忍耐而合谋，为忍耐而结盟；他们在对德行的虚饰中一致追求忍耐；他们关于忍耐展示出一切智慧之夸耀。这是对忍耐的伟大见证，因为它甚至激发世间虚浮的学派前来赞美和荣耀它！或者，这反倒是一种损害，因为神圣之物竟在世俗学识中流连？但那些不久将因自己的智慧而感到羞耻（这智慧将随它所寄居的世界一同毁灭和蒙羞）的人，且让他们自己去面对此事吧。\par

% structure: p0004 source=h2 target=omitted reason=duplicate-page-title

对于\Emph{我们}，没有哪种人为的、由麻木塑造的犬儒式平静，为操练忍耐提供依据；而是那活生生的、属天的训导之神圣安排，首先将天主本身摆在我们面前作为忍耐的榜样。祂将这光芒的光辉均匀地散布给义人和不义之人；祂容忍季节的恩惠、元素的服务、整个自然的贡赋同样临于堪当者和不堪当者；祂容忍最忘恩负义的民族（他们崇拜技艺的玩物和双手的作品，迫害祂的圣名连同祂的家人）；祂容忍奢侈、贪婪、不义、邪恶，日益狂妄自大；以致祂以自己的忍耐贬损了自己；因为许多人之所以不信主，正是因为他们长久以来不知道祂对世界发怒。\par

% structure: p0006 source=h2 target=section reason=relative-heading-rank; base=h2

\section{\Emph{第三章 耶稣基督在降生成人及其工程中提供了更可效法的忍耐榜样}}

然而，这种天主性的忍耐，仿佛远在天边，或许被视为\Quote{为我们太高的事}；但那种在人间公开被人亲手触摸到的\BibleRef{若一 1:1}忍耐是什么呢？天主容许自己受孕于母腹，等待\Emph{诞生的时期}；出生后，忍受\Emph{成长的延迟}；长大成人后，并不急于被人认出，反而羞辱自己，由自己的仆人施洗；只用言语击退诱惑者的攻击；从\Quote{主}变为\Quote{师傅}，教导人逃避死亡，而他自己却已受训于被冒犯的忍耐之绝对的宽容。他不争竞，不喧嚷，街上也无人听见他的声音。压伤的芦苇他不折断，将残的灯火他不吹灭：因为先知——不，是天主自己的见证，将他自己的圣神连同完全的忍耐一同放在他儿子里面——并没有说谎。没有哪个渴望依附他的人是他所不接纳的。没有一人的桌子或屋顶是他所鄙视的：他自己甚至为门徒洗脚；他不拒绝罪人，也不拒绝税吏；他甚至对那拒绝接待他的城也不发怒\BibleRef{路 9:51}，虽然门徒曾希望从天上降下烈火立刻焚烧那如此无礼的城。他照顾忘恩负义的人，他向陷害他的人让步。这还算小事，如果他不是甚至容忍自己的出卖者在身边，并坚决不指认他。不仅如此，当他被出卖、被牵去\BibleQuote{如同羊被牵去宰杀}（因为\BibleQuote{他在剪毛者手下默不作声，如羔羊一般}）时，他本可以一句话就召来天上成营的天使，却连一个门徒的复仇之剑也不认可。主的忍耐因\Person{玛耳曷}的伤口而受伤。他也诅咒将来一切刀剑的行为；并通过恢复健康，使那他没有伤害的人得到满足——这都是由于忍耐，仁慈之母。我沉默地略过他被钉十字架的事实，因为他正是为此而来；然而，那必须承受的死亡，难道也需要同样多的凌辱吗？\Emph{不}，在他将要离去时，他渴望饱享忍耐的喜悦。他被吐唾沫、被鞭打、被嘲笑、被粗鄙地穿上衣服，更粗鄙地戴上冠冕。镇静的信仰何等奇妙！那为了隐藏自己而取了人形的主，丝毫没有模仿人的不耐！因此，\Person{法利塞}人，你们本应比其他任何特征更早认出主来。\Emph{人}中没有一个能成就这种忍耐。如此伟大而有力的证据——其巨大程度在外邦人中确实是信仰被拒绝的原因，但在我们中却是信仰的理由和根基——已足够清楚地表明（不仅在命令的讲道中，同样也在主忍受的苦难中），对于蒙赐相信的人而言，忍耐作为某种固有德性的结果和卓越，乃是天主的本性。\par

% structure: p0008 source=h2 target=section reason=relative-heading-rank; base=h2

\section{第四章. 效法我们师傅的义务——由奴隶甚至野兽教导我们。顺服的效仿建立在忍耐上}

因此，如果我们看见所有正直合宜的仆人都按照主人的性情塑造自己的行为；也就是说，得宠的艺术是服从，而服从的规则是顺从的归属：那么，我们岂不更应当显出一种与我们的主相称的品格吗？——我们是永生天主的仆人，他的审判不取决于锁链或自由帽，而取决于惩罚或救恩的永恒；要逃避这严厉或赢得这慷慨，就需要一种与严厉所发出的威吓或慷慨所自由赐予的应许同样巨大的殷勤服从。然而，我们不仅从\Emph{人}那里要求服从——他们颈下有奴隶的束缚，或以其他法律方式欠下服从的债——甚至从牲畜、从野兽那里，我们也要求服从；我们明白它们是主为我们提供和交付的。那么，天主使我们服于它们的\Emph{受造物}，在服从的纪律上比我们更好吗？最后，那些服从的\Emph{受造物}认识\Emph{它们的主人}。我们是否犹豫要仔细听从那唯一我们所属的主——即主？但是，不从你自己身上回报那同样的事——即你通过你邻舍的宽容从别人那里得到，却回报给那使你得到的人——是多么不义，多么忘恩啊！关于我们必须向主天主所表现的服从，无需多言；因为对天主的认识懂得它当尽的责任。然而，为了不使我们关于服从的\Emph{言论}显得不相干，（让我们记住）服从本身来自忍耐。一个\Emph{不耐}的人绝不会这样做，而一个\Emph{忍耐}的人必以此为乐。那么，谁能充分论述那\Emph{忍耐}的善呢？主天主，一切善事之显示者与接纳者，曾在自身中带着这忍耐。再者，谁又会怀疑，每一件善事，因为它属于天主，就应当被属于天主的人以全心热烈追求呢？藉着这些（考量），对忍耐的称许与劝勉被简要地、几乎是在一条规训规则的纲要中建立起来了。\par

% structure: p0010 source=h2 target=omitted reason=duplicate-page-title

然而，对信仰必需之事进行讨论的过程并非徒劳，因为它并非不结果实。在建立方面，没有喋喋不休是卑劣的，如果它有时是卑劣的话。所以，如果话语是关于某种特定的\Emph{善}，主题要求我们也考察那善的\Emph{相反}。因为如果你首先总结了什么应当避免，你就会更清楚地阐明什么应当追求。\par

因此，让我们思考\Emph{不耐}这一品质——正如忍耐存在于天主内一样，它的敌对品质是否也在我们的仇敌身上产生并被察觉，好由此看出它如何从根本上与信德相悖。因为由天主的对手所孕育的，当然不会与天主的事物友善。\Emph{事物}的不和也就是它们\Emph{作者}的不和。再者，既然天主是至善的，魔鬼则相反是\Emph{存有中最恶的}，两者因自身的极端差异已证明彼此不为对方工作；因此善者不能为我们行任何恶事，恶者也不能行任何善事。所以，我在魔鬼自身中发现了不耐的诞生，就在他无法忍受主天主将祂所造的万物都交托给祂自己的肖像，即人的时候。因为他若忍受了，就不会悲伤；若不悲伤，就不会嫉妒人。于是，他因嫉妒而欺骗了人；他嫉妒是因为他悲伤；他悲伤则是因为他没有忍耐。这个堕落的元始天使究竟是先有了恶意还是先有了不耐——我不屑于追问；因为显而易见，要么不耐\Emph{与恶意一同}兴起，要么恶意\Emph{源于}不耐；随后两者相互勾结，在同一个父的怀抱中不可分割地成长。然而，他凭自己的经验得知，自己所首先感受并借此走上犯罪道路的，对犯罪是何等有力的帮助，于是他便用它来诱使人陷入罪恶。女人一遇到他——我这样说并不冒昧——通过他与她的对话，就被一种染有不耐的精神所吹拂；如此肯定的是，她若始终持守忍耐以尊重天主的命令，就根本不会犯罪。她怎能忍受并非独自遇见了魔鬼？在亚当面前，那时他尚不是她的丈夫，也无需倾听她的话，她却不耐烦保持沉默，使他成为她从恶者所领受之事的传递者？于是，另一个人也因一人的不耐而丧亡；随即，他自己也因自己的不耐在两方面——既对天主的预告，也对魔鬼的欺骗——而丧亡，既不能忍受遵守前者，也不能忍受驳斥后者。因此，犯罪的开端从此而来，审判的最初起源也从此而来；人从此被诱导冒犯天主，天主开始发怒。天主的第一次义怒由此而来，祂的第一次忍耐也由此而来；祂那时只以诅咒为满足，在魔鬼的事上克制了立即施罚。否则，在这不耐心之罪以前，人还有什么罪可归咎呢？他是无辜的，与天主亲密无间，是乐园的农夫。但一旦他屈服于不耐，就完全失去了在天主前的馨香；完全无法承受天上的事物。从此，一个归于土、被逐出天主视线的受造物，开始因不耐轻易转向一切冒犯天主的行为。因为那由魔鬼种子孕育的\Emph{不耐}，在恶的丰饶中产生了愤怒作为她的儿子；而一旦生出，她便用自己的技艺训练他。因为那使亚当和厄娃沉入死亡的事，也教导了他们的儿子以杀人为开端。若加音——第一个杀人和第一个杀兄弟的人——能平和忍受（而非不耐烦地）主拒绝他的祭品，若他不向自己的兄弟发怒，若他最终没有夺取任何人的生命，那我将徒劳地将其归于不耐。既然他若非发怒就不会杀人，若非不耐就不会发怒，他便表明他因愤怒所行的事必须归因于那在不耐的摇篮时期——当时尚在襁褓中——刺激愤怒产生的东西。但后来它的增长何其巨大！不足为奇，若它是最初的犯罪者，那么因为它\Emph{就是}最初的，因此它便也是\Emph{一切}犯罪唯一的母根，从自己的源头倾泻出各种罪恶的支流。我们已经谈到\Emph{谋杀}；但因为它从最初就是\Emph{愤怒}的结果，随后它为自己找到的其他任何原因，都一并归咎于不耐，作为其根源。因为无论出于\Emph{私仇}还是为了\Emph{贪利}而行此恶，第一步都是他变得\Emph{不耐}于仇恨或贪财。任何\Emph{迫使}一个人的事，若没有\Emph{对自身的不耐}，不可能在\Emph{行为}中完成。谁\Emph{曾}在\Emph{不耐情欲}的情况下犯\Emph{奸淫}？再者，若女性因价格而被迫出卖贞操，当然是由不耐轻视利益而支配这\Emph{出卖}。我提及这些作为主眼中主要的罪行，因为简言之，\Emph{每一项}罪都可归因于不耐。\Quote{恶}就是\Quote{对善的不耐}。没有\Emph{不贞}的人不是\Emph{对贞洁不耐}；\Emph{不诚实}的人对\Emph{诚实不耐}；\Emph{不敬}的人对\Emph{虔敬不耐}；\Emph{不安}的人对\Emph{安宁不耐}。为使每个人成为\Emph{恶}，他将\Emph{不能坚忍}于\Emph{善}。因此，这样一个罪行九头蛇怎能不冒犯主——恶的谴责者？难道不明显，连以色列自己也是因不耐而总是在对天主的义务上跌倒，从那时起——当他忘记了曾将他从埃及苦难中拉出的上天手臂，向亚郎要求\Quote{众神作他的引导}；他将自己的金子贡献给偶像：为了梅瑟遇见天主所需的必要耽搁，他却不耐烦忍受。在可食的玛纳之雨和岩石流出的水之后，他们因不能忍受三天的干渴而对主绝望；这也被主归咎为他们不耐。并且——不要逐例漫游——没有哪次他们不是因不耐而失职丧亡。再者，他们如何下手对付先知？除非因不耐听他们。再者，他们如何对待主自己？同样因不耐看见祂。但他们若走上\Emph{忍耐的道路}，就会获得解救。\par

% structure: p0013 source=h2 target=section reason=relative-heading-rank; base=h2

\section{第六章 忍耐在信德之前与之后}

因此，忍耐既是信德之后，也是信德之前。简言之，亚巴郎信了天主，被祂视为正义；但正是忍耐证明了他的信德，当他奉命祭献自己的儿子时——我并非要说试探，而是要说对他信德的典型印证。但天主知道祂所视为正义的是谁。如此重大的命令，即使完全执行也不蒙主喜悦，他却忍耐地听了，并且（若天主愿意）也必会执行。因此，他被称为\Quote{有福的}，因为他\Quote{忠信}；他被称为\Quote{忠信的}，因为他\Quote{忍耐}。这样，信德被忍耐光照，当它藉着\BibleQuote{亚巴郎的后裔，就是基督}\BibleRef{迦 3:16}在万民中传播，并将恩宠加于法律之上时，便使忍耐成为她卓越的助手，以扩增并满全法律，因为这是正义教导中所唯一欠缺的。因为古人惯常要求\Quote{以眼还眼，以牙还牙}，并以利息偿还\Quote{以恶报恶}；因为那时忍耐尚不在世上，因为信德也不在。当然，在此期间，不耐烦享受着法律所给予的机会。当忍耐的主人和师傅缺席时，这是容易的。但祂降临后，将信德的恩宠与忍耐结合，如今连用言语攻击，甚至说\Quote{愚笨}的话，也不可以，否则就有\Quote{受审判的危险}。愤怒已被禁止，我们的心神被保留，手的粗暴被制止，舌头的毒汁被拔出。法律所得的比所失的更多，因为基督说：\BibleQuote{爱你们的仇敌，为诅咒你们的人祝福，为迫害你们的人祈祷，好使你们成为你们在天之父的子女。}\BibleRef{玛 5:44}你看见忍耐为我们赢得了怎样的父亲吗？在这首要的诫命中，忍耐的全盘操练被简洁地包含，因为即使当恶行是应得的，也不再被容许。\par

% structure: p0015 source=h2 target=section reason=relative-heading-rank; base=h2

\section{第七章 不耐烦的原因及其相应的诫命}

如今，当我们遍历不耐烦的原因时，所有其他的诫命也将在各自的位置上回应。若我们的心神因财产的损失而被激起，主在圣经中几乎处处都劝勉我们轻视世界；也没有比主本人不置身于任何财富之中这一事实给予我们更强大的劝勉来轻视金钱。祂总是为穷人辩护，预先谴责富人。因此，祂预先为忍耐准备了\Quote{损失}，为富足准备了\Quote{轻视}（作为分份）；藉着（祂自己）对财富的\Emph{弃绝}，表明对财富的损害也不应太在意。因此，对于那根本不需要追求的东西——因为主也没有追求它——我们应当忍受其削减或剥夺而不心痛。主的神曾藉着宗徒宣称\Quote{贪财}是\Quote{一切邪恶的根源}。我们不要认为那\Quote{贪财}仅仅在于贪图他人的东西：因为甚至那些\Emph{似乎}是我们的东西\Emph{也是}他人的；因为没有什么属于我们，一切都是天主的，我们也属于祂。所以，若在遭受损失时，我们感到不耐烦，为失去并非自己的东西而悲伤，我们就暴露出接近贪财的倾向：当我们难以忍受失去他人的东西时，我们是在\Emph{寻求}他人的东西。因损失而极度激动不耐烦的人，因将地上的事物置于天上的事物之前，直接得罪了天主；因为他为了世上的事物大大震动了他从主所领受的圣神。因此，让我们乐意失去地上的事物，持守天上的事物。让整个世界消亡，只要我能获得忍耐！事实上，我不知道那个下定决心不坚忍地承受某物损失的人（无论损失是由于偷窃、强暴，甚至疏忽），会否乐意并真心地为了施舍而动用他自己的财产？因为那完全不能忍受别人切割的人，自己却在身体上拔剑？损失中的忍耐是施予和分享的操练。不怕失去的人，不会厌于施予。否则，当一个人有两件外衣时，他怎能将一件给赤身露体的人\BibleRef{路 3:11}？除非他也是那种连外衣也被拿走时愿意连衬衫也给人的人？我们怎能藉着\OriginalTerm{玛门}{Mammon}为自己结交朋友\BibleRef{路 16:9}，若我们如此爱它以至于不能忍受失去它？我们将与失去的\OriginalTerm{玛门}{Mammon}一同灭亡。为何我们在这里\Emph{找到}，而我们的任务却是\Emph{失去}？对一切损失表现不耐烦是外邦人的事，他们或许把钱置于灵魂之上；因为他们正是这样做的：当他们贪图利润而遭遇海上贸易的危险时；为了钱，即使在市场上，也没有任何他们不敢尝试的连地狱也惧怕的事；当他们为娱乐和军营而雇用时；当他们像野兽一样在公路上打劫时。但我们，根据我们与他们区别的不同，应当不是为了钱而舍弃灵魂，而是为了灵魂而舍弃钱，无论是自愿施予还是忍耐地损失。\par

% structure: p0017 source=h2 target=omitted reason=duplicate-page-title

我们这些在此世界上使自己的灵魂和身体暴露于各种伤害之下、并在伤害中展现忍耐的人，难道会在失去次要事物时感到受伤吗？愿基督的仆人远离这样的玷污，即那已为更大诱惑预备好的忍耐竟在琐碎之事上离弃他。若有人试图用暴力打你，主的训诫就在眼前：祂说：\Quote{那打你脸的，连另一边也转给他。}\BibleRef{玛 5:39}让暴怒被你的忍耐耗尽。无论那打击如何，伴随着痛苦和侮辱，它必将在主那里受到更重的打击。你以忍耐更伤害那暴怒者：因为那为你忍耐的缘故，主将击打他。若舌头的苦毒以咒骂或责备爆发，请回顾这句话：\Quote{当他们咒骂你们时，你们要欢喜。}主自己在法律眼中被咒骂，而祂却是唯一蒙福的那位。因此，让我们这些仆人紧紧跟随主，并忍耐地被咒骂，以便我们能够蒙福。若我以过于平静的心态听到针对我的轻浮或邪恶的话语，我必定要么自己报复那苦毒，要么被无声的不耐烦折磨。那么，当我被咒骂时，我（用舌头）反击，怎能算作跟随了主的教导？这教导说：\Quote{人不是因器皿的污秽而污秽，而是因从他口中发出的事物而污秽。}此外，经上说：\Quote{每一句虚妄和无益的话都将使我们受控告。}由此可知，主禁止我们做什么，祂就劝勉我们以忍耐从他人那里承受什么。我还要加上一点关于忍耐的\Emph{乐趣}。因为每一次伤害，无论来自舌头或手，一旦落在忍耐上，就会如同某种武器击中最坚硬的磐石而被挫钝一样被消除。因为它将立即在徒劳无功的努力中坠落；有时它会反弹，以反击的势头向发出者发泄怒气。无疑，有人伤害你，是为了使你痛苦；因为伤害者的快乐在于受害者的痛苦。因此，当你因不痛苦而粉碎他的快乐时，\Emph{他}必然因失去快乐而痛苦。于是你不仅毫发无损地离开——这本身已足够——而且还因对手的失望而满足，因他的痛苦而报复。这就是忍耐的\Emph{益处}和\Emph{乐趣}。\par

% structure: p0019 source=h2 target=omitted reason=duplicate-page-title

就连失去亲人时的那种不耐烦，也不可原谅——那种以对悲伤的权利为庇护的辩词。因为我们必须将宗徒的宣言摆在面前，他说：\Quote{不要因任何人安眠而过度悲伤，如同那些没有希望的外邦人一样。}这是合理的；因为相信基督的复活，我们也相信自己的复活，基督正是为我们而死、并为我们而复活。既然死者的复活是确定的，为死亡而悲伤便是多余的，悲伤的不耐烦也是多余的。因为如果你相信（你所爱的人）并未消亡，为何还要悲伤？为何要焦急地承受他的暂时离开，而你相信他会回来？你认为的死亡，其实是离去。那先走一步的人不应被哀悼，尽管无疑应被思念。那种思念也须以忍耐调和。因为为何要毫无节制地忍受\Emph{他离去}的事实，而你不久就会跟随？此外，这类事情上的不耐烦对我们的希望是不祥之兆，是对信仰的不真诚。当我们不以平静之心接受主从这世界召唤任何人，仿佛他们值得怜悯时，我们就伤害了基督。宗徒说：\Quote{我渴望现在就被接去，与基督同在。}他所展示的渴望是多么美好啊！那么，如果我们对那些已经获得基督徒渴望的人不耐烦地悲伤，我们自己便显示出不愿获得这渴望。\par

% structure: p0021 source=h2 target=section reason=relative-heading-rank; base=h2

\section{第十章 论报复}

还有另一种不耐烦的主要刺激，即报复的私欲，它或是关乎荣显之事，或是关乎恶意。但是，一方面，\Quote{光荣}处处是\Quote{虚幻的}；另一方面，恶意在主眼中总是可憎的；在此情形下尤其如此，因为当它被邻人的恶意所激怒时，它自命在实施报复上更为优越，并以恶报恶，使已行的事加倍。在错误的判断中，报复似乎是痛苦的慰藉；相反，在真理的判断中，它被证实为恶毒。因为挑事者与被激怒者之间有何区别，除了前者被发现在行恶上是先者，而后者是后者？然而，二人在主眼中都被指控伤害人，主禁止并谴责一切恶行。在行恶上，不考虑\Emph{次序}，也不因\Emph{场合}而分开\Emph{相似性}所结合的事。并且诫命是绝对的：不可以恶报恶。\BibleRef{罗 12:17} 同样行为导致同样功过。我们如何遵守那原则，如果我们憎恶时却不憎恶\Emph{报复}呢？此外，我们若将报复的仲裁权归于自己，将向主天主献上什么尊敬呢？我们是腐败的——\BibleRef{依 64:6}——瓦器。对于我们的仆役，如果他们擅自对同伴施行报复，我们会非常恼怒；而那些以忍耐奉献给我们的人，我们不仅赞赏他们记念谦卑、仆役身份，并热切关切主尊荣的权利，而且我们给他们比他们自己所预先要求的更充分的满足。对于如此公正判断、如此有力执行的主，难道会有不同结果的风险吗？那么，我们为何相信祂是审判者，如果祂不同时是报复者呢？祂应许要为我们这样做，说：\Quote{报复\Emph{属于}我，我必要报复。}也就是说：\Emph{将}忍耐留给我，我将酬报忍耐。因为当祂说：\Quote{不要判断，免得你们受判断}，祂不是在要求忍耐吗？谁若不在报复自己上忍耐，谁又能克制不判断别人呢？谁\Emph{判断}是为了\Emph{宽恕}？如果他宽恕，他仍然顾念了判断者的不耐烦，并夺走了那唯一判断者，即天主的尊荣。这种不耐烦曾遭遇多少不幸！它多少次后悔自己的报复！它的狂暴多少次比导致它的原因更恶劣！——因为凡以不耐烦承担的事，无不带着冲动完成；凡带着冲动做的事，无不跌倒，或全然坠落，或一头栽下去。此外，如果你报复得太轻，你会发狂；如果报复得太重，你将背负重担。我与报复有何干系，我因痛苦的不耐烦而无法调节其程度？然而，如果我靠忍耐安息，我将不会\Emph{感到}痛苦；如果我不感到痛苦，我将不会\Emph{渴望}报复自己。\par

% structure: p0023 source=h2 target=section reason=relative-heading-rank; base=h2

\section{第十一章 实践忍耐的进一步理由。它与真福的关系}

在尽我们所能列举了这些不耐烦的主要物质原因之后，我们为何要绕路去考虑其余的原因——有哪些是在家中的，有哪些是在外的？恶者的运作广泛而扩散，投掷我们心灵的多种刺激，有时是琐碎的，有时是非常大的。但琐碎的你可以因其微小而轻视；非常大的你可以因其压倒性而屈服。伤害越小，就越不需要不耐烦；伤害越大，就越需要对付伤害的补救措施——忍耐。因此，让我们努力承受恶者的打击，好使我们平静的忍耐嘲笑敌人的狂热。然而，如果我们自己，或因不谨慎，或主动地，招致任何事，让我们以同样的忍耐来承受我们不得不责备自己的事。此外，如果我们相信某些打击是主加给我们的，我们更应向谁显出忍耐，难道不是向主吗？不仅如此，祂还教导我们感恩和喜乐，因为被认为配得上天主的管教。祂说：\Quote{我所爱的，我必管教。}哦，有福的仆人，主关注他的改进！祂屈尊向他发怒！祂不掩饰自己的责备以欺骗他！因此，我们处处都有责任履行忍耐的义务，无论我们从何方，或因自己的错误，或因恶者的陷阱，招致主的责备。那义务的赏报是大的——即幸福。因为主称谁为有福的呢？除了忍耐的人？祂说：\Quote{神贫的人是有福的，因为天国是他们的。}\BibleRef{玛 5:3} 确实，除非是谦卑的人，没有人是\Quote{神贫的}。那么，谁是谦卑的呢？除非是忍耐的人？因为没有人能在不首先忍耐承受自卑行为的情况下自卑。祂说：\Quote{哀恸和哭泣的人是有福的。}\BibleRef{玛 5:4} 因此，向这样的人应许了\Quote{安慰}和\Quote{欢笑}。\Quote{温良的人是有福的：}\BibleRef{玛 5:5} 在这个术语下，不耐烦的人肯定不可能被归类。同样，当祂以同样的福乐称号标记\Quote{缔造和平的人}，\BibleRef{玛 5:9} 并称他们为\Quote{天主的子女}时，请问不耐烦的人与\Quote{和平}有任何相似之处吗？即使是愚人也看得出。然而，当祂说：\Quote{几时他们辱骂并迫害你们，你们要欢喜踊跃，因为你们在天上的赏报是大的。}当然，祂并非向踊跃的\Emph{不耐烦}应许那赏报；因为没有人\Emph{会}在逆境中\Quote{踊跃}，除非他首先学会轻视它们；除非学会实践忍耐，否则没有人会轻视它们。\par

% structure: p0025 source=h2 target=section reason=relative-heading-rank; base=h2

\section{第十二章 某些其他神圣诫命。宗徒的爱德描述。它们与忍耐的关系}

至于和平的规条，它如此中悦天主；在这个倾向于不耐的世界里，有谁会甚至\Emph{一次}就宽恕他的弟兄？我不说\Quote{七次}，或\Quote{七十七次}？有谁思量要与对头对簿公堂，能以协议了结此事\BibleRef{玛 5:25}，除非他首先砍掉烦恼、硬心和苦毒——这些实际上是不耐的有毒产物？你将如何\BibleQuote{赦免，也得蒙赦免}\BibleRef{路 6:37}，若缺少忍耐使你固执于不义？凡是心中与弟兄不和睦的人，他无法完成在祭台上献\BibleQuote{虔诚的礼物}，除非他先存心\BibleQuote{与弟兄和好}而回到忍耐。\BibleRef{玛 5:23}若\Quote{太阳落在我们的忿怒上}，我们就处于危险中：我们不容许没有忍耐地度过一天。然而，既然忍耐在每一种有益的纪律中居于首位，那么她同样服务于悔改（悔改习惯于搭救那些跌倒的人），有什么可奇？当婚姻关系解除时（我指那种使丈夫或妻子得以持守永久守寡的合法原因），她（忍耐）等待、切望、用恳求说服所有终将进入救恩之人的悔改。她赐予每个人多大的祝福！她阻止一人成为奸夫；她改正另一人。同样，她在主比喻中关于忍耐的那些神圣榜样中出现。牧人的忍耐寻找并找回了迷路的母羊\BibleRef{路 15:3}：因为\Emph{不}耐会轻易鄙视\Emph{一只}母羊；但忍耐承担了寻找的劳苦，那耐心的负重者将遗弃的罪人扛在肩上带回家。那个荡子，父亲的忍耐也接纳他，给他穿衣、喂食，并在愤怒的兄长的\Emph{不}耐面前为他辩护\BibleRef{路 15:11}。因此，那\Quote{丧亡了的}得救了，因为他踏上了\Emph{悔改之}路。悔改不会丧亡，因为它找到了忍耐（来迎接它）。因为爱德——信德的最高圣事，基督徒名号的宝库，宗徒以圣神的全部力量所赞美的——除了忍耐的教训之外，还能由谁的教训而受训练呢？他（宗徒）说：\BibleQuote{爱德是含忍的；}如此她应用了忍耐：\BibleQuote{是慈祥的；}忍耐不做恶事：\BibleQuote{不嫉妒；}那无疑是忍耐的特殊标记：\BibleQuote{不夸张暴力；}她将自己的节制从忍耐中汲取：\BibleQuote{不自大；不粗暴；}因为那不属于忍耐：\BibleQuote{不追求自己的利益；}如果她\Emph{付出}自己的，只要能使邻人受益：\BibleQuote{不轻易发怒；}若她发怒，她还能给\Emph{不}耐留下什么？因此他说：\BibleQuote{爱德凡事包容，凡事忍耐；}当然因为她忍耐。因此，正当地，\BibleQuote{她永不失败；}因为其他一切都要取消，将要结束。\BibleQuote{语言、知识、先知之恩，都将停止；信德、望德、爱德，永存：}信德，由基督的忍耐引入；望德，由人的忍耐等候；爱德，由忍耐伴随，以天主为师。\par

% structure: p0027 source=h2 target=section reason=relative-heading-rank; base=h2

\section{第十三章 论身体的忍耐}

至此，最后，论及纯一不变的忍耐，以及它纯粹存在于\Emph{心灵}中的情形：但我同样以多种方式在\Emph{身体}上追求它，目的在于\Quote{赢得主}；\BibleRef{斐 3:8}因为这也是主亲自在身体的美德上所彰显的一种品质；如果统治理智确实易于将其圣神的恩赐与它的\Emph{身体}住所沟通的话。那么，忍耐在\Emph{身体}上的事务是什么呢？首先，\Emph{它是}肉身的磨难——一个能通过谦卑的祭献平息主的牺牲品——通过向主奠祭卑贱的衣袍，加上食物的匮乏，满足于简单的饮食和纯净的水，同时将禁食\Emph{与这一切}结合起来；使自己习惯于苦衣和灰烬。这种\Emph{身体}上的忍耐为我们的善祷增添恩宠，为我们的恶祷增强力量；这开启了基督\Emph{我们的}天主的耳朵，驱散严厉，引出仁慈。因此，那位巴比伦王，在七年的污秽和忽视中被逐出人形，因为他得罪了主；通过身体上的忍耐祭献，不仅恢复了王国，而且——这是人更渴望的——向天主做了补赎。此外，如果我们按顺序列出更高、更幸福的身体上的忍耐的等级，（我们发现）正是她被圣洁委托掌管肉身的节制：她保守寡妇，给贞女盖上印记\BibleRef{格前 7:34}，并将自阉的阉人提升到天国。\BibleRef{玛 19:12}那源于\Emph{心灵}德行的东西在\Emph{肉身}中被成全；最后，通过肉身的忍耐，在迫害下作战。如果逃跑紧迫，肉身与逃跑的不便作战；如果监禁降临我们，肉身（仍然）被捆绑，肉身在镣铐中，肉身在独居中，在那缺少光明中，在那忍受世界虐待的忍耐中。然而，当它被引向幸福的最终考验、引向第二次圣洗的机会\BibleRef{路 12:50}、引向登上神圣宝座的行动时，没有比\Emph{身体}上的忍耐更需要\Emph{的}了。如果\Quote{精神固然愿意，但肉身，} \Emph{没有}忍耐，\Quote{软弱，} \BibleRef{玛 26:41}那么，除了\Emph{在忍耐中}，精神和肉身本身的安全何在？但是当主论及肉身这样说话，宣布它\Quote{软弱，}时，祂显示需要加强它——即通过忍耐——以应对所有颠覆或惩罚信仰的准备；使它恒心忍受鞭打、火焰、十字架、野兽、刀剑；所有这一切，先知和宗徒们通过忍受而征服了！\par

% structure: p0029 source=h2 target=section reason=relative-heading-rank; base=h2

\section{第十四章 这双重忍耐——灵性上的与身体上的——的能力。在古圣先贤身上的例证}

凭着这忍耐的力量，依撒意亚被\Emph{锯成两半}，却不停息地谈论主；斯德望被石头砸死，却为他的敌人祈求宽恕。\BibleRef{宗 7:59}哦，他也真有福，他凭借各种忍耐的努力，承受了魔鬼的一切暴行！——无论是牛群的被夺，还是羊群的财富，还是子女在一场灾祸中全部被扫除，最后，还是他自己身体（全身）伤痛的折磨，都没有使他疏远他向主所誓许的忍耐和信德；魔鬼全力击打他，却徒劳无功。因为他所有的痛苦都没有使他远离对天主的敬畏；相反，他被立为我们的榜样和见证，以彻底完成既在精神上也在肉身中，既在心灵上也在身体上的忍耐；好使我们不屈服于世俗财物的损失，不屈服于最亲爱之人的丧失，甚至不屈服于身体的痛苦。天主在那位英雄身上为魔鬼设立了何等样的灵柩！祂为祂荣耀的敌人竖立了何等样的旌旗，当那人每次听到苦涩的消息时，嘴里只发出对天主的感谢，而斥责他的妻子，她已被诸般痛苦折磨得疲惫不堪，并催促他采取邪僻的补救办法！天主如何微笑，恶者如何被割裂，当约伯以极大的镇定不断地刮去\BibleRef{约 2:8}自己溃疡的不洁溢出物，并戏谑地将从那里爆发出来的寄生虫放回他斑驳肉身的同一洞穴和饲料场！因此，当所有诱惑的箭矢都因他的忍耐之胸甲和盾牌而变钝时，那天主胜利的工具不仅立即从天主那里恢复了身体的健康，而且加倍地拥有了他所失去的一切。如果他愿意他的子女也复活，他可能会再次被称为父亲；但他宁愿他们\Quote{在那一天}复活。那种喜乐——完全确信主——他推迟了；同时他忍受了一种自愿的丧子之痛，好使他不会没有某种（操练）忍耐而生活。\par

% structure: p0031 source=h2 target=omitted reason=duplicate-page-title

天主实在是忍耐之无比充足的寄托者。若您将委屈交托于祂，祂是复仇者；若损失，祂是恢复者；若痛苦，祂是医治者；若死亡，祂是复活者。忍耐何等荣幸，竟使天主成为她的债务人！并非无因：因她持守祂的一切法令，执行祂的一切诫命。她坚固信德；是和平的舵手；扶助爱德；建立谦逊；长久等待悔改；为告解盖上印记；管束肉身；保存心神；勒住舌头；约束手；将诱惑践踏在脚下；驱散绊脚石；为殉道者加冕恩宠；安慰穷人；教导富人节制；不使弱者过度劳累；不耗尽强者；是信友的喜乐；吸引外邦人；将仆人荐于主人，将主人荐于天主；装饰妇女；使男子蒙悦纳；童年受人喜爱，青年受人赞美，老年受人敬仰；无论男女，一生各阶段，皆显美丽。来，现在\Emph{看}我们是否对她的仪容和习性有了大致概念。她的面容平静安详，眉宇宁静，无愁无怒之纹；双眉舒展愉悦，双目谦卑低垂，非因忧伤；口唇以沉默的高贵印记封缄；面色如无忧无咎之人；频频摇头对抗魔鬼，笑中含威；此外，她胸前衣袍洁白合身，既非膨胀亦非凌乱。因为\Emph{忍耐}坐在那最温和最温柔的圣神的宝座上；祂不在旋风滚滚中，不在乌云铅色里，而是柔和宁静，坦率纯朴，即\Person{厄里亚}在第三次尝试时所见的。因为天主在哪里，祂的养子——忍耐——也在那里。当天主的圣神降临时，忍耐便紧随祂不可分离。若我们和圣神同来时不让她进入，祂会永远停留吗？不，我不知道祂是否还会\Emph{再}停留。没有这位同伴和侍女，祂必然处处时时受窘。无论祂的仇敌施加什么打击，祂都无法独自承受，因为没有承受的工具手段。\par

% structure: p0033 source=h2 target=section reason=relative-heading-rank; base=h2

\section{第十六章 外邦人的忍耐与基督信徒的忍耐迥然不同。他们的忍耐注定丧亡，我们的忍耐注定得救。}

这是规则，这是纪律，这是忍耐的作为，\Emph{即}天上的、真正的忍耐；也就是基督信徒的\Emph{忍耐}，而非虚假可耻的，如同地上万民的忍耐一样。因为魔鬼为了在这一点上也与主竞争，他几乎完全对等地（只是邪恶与善良的多样性恰好与它们的幅度相等）也教导了\Emph{他的门徒}一种他特有的忍耐；我是指那种使丈夫因嫁妆而变得唯利是图，又教他们从事拉皮条的交易，使他们受制于妻子的权力；那种以虚伪的情感，承受一切勉强奉承的劳苦，旨在诱骗无子嗣者；那种使腹欲的奴仆屈从于侮辱性的庇护，将自由置于咽喉之下。外邦人熟悉这样的忍耐行径；他们急切地用如此美善的名号来掩盖丑恶行径：他们忍耐地对待竞争对手、富人以及宴请他们的人；唯独对天主不能忍耐。但让他们自己和他们的首领的忍耐自己去面对吧——地下之火正等着\Emph{那种忍耐}！至于我们，要热爱天主的忍耐、基督的忍耐；我们要将祂为我们付出的\Emph{忍耐}回报给祂！我们要将心神上的忍耐、肉体上的忍耐奉献\Emph{给祂}，因为我们相信肉身和灵魂的复活。\par

\SourceRecord{Translated by S. Thelwall.；Ante-Nicene Fathers；Vol. 3.；Edited by Alexander Roberts, James Donaldson, and A. Cleveland Coxe.；Buffalo, NY: Christian Literature Publishing Co.,；1885.；<http://www.newadvent.org/fathers/0325.htm>.}
